% Chapter 1: Introduction.

\section{Background of the Study}
Stock exchanges and regulators watch for trading days that depart from normal market
behaviour, such as sudden crashes, manipulated prices or reactions to regulatory news.
Finding such days automatically is an anomaly detection problem, and deep learning models
that learn normal behaviour from data and score departures from it are now widely used for
time series~\cite{malhotra2016lstm,hundman2018detecting,xu2018unsupervised}.

Financial returns are a difficult setting for these models. Their variance is not constant:
large moves cluster in turbulent periods and small moves in calm ones, a property captured
by ARCH and GARCH models~\cite{engle1982autoregressive,bollerslev1986generalized}. A 3\%
fall is ordinary in a crisis and exceptional in a quiet year. A detector that judges every
day on the same scale will mark crisis periods as anomalous and miss unusual days in calm
periods.

The difficulty is greater in emerging and frontier markets. The Dhaka Stock Exchange (DSE)
in Bangladesh has a broad index, DSEX, only from January 2013. Its prices are not adjusted
for the frequent bonus share issues of listed firms, and for long stretches between 2020 and
2024 the Bangladesh Securities and Exchange Commission (BSEC) imposed floor prices that froze
many stocks. Little clean data is available to train a detector on the DSE alone.

\section{Problem Statement}
Published DSE volatility studies fit GARCH-family models in
sample~\cite{siddikee2016volatility,haque2021modeling}, and published anomaly detection
studies on stock data do not adjust their inputs for volatility
clustering~\cite{golmohammadi2015time,mariani2021autoencoder}. Among the work reviewed in
Chapter~2, none applies a deep anomaly detector to the DSE, none uses GARCH-standardized
returns as a detector's input, and none transfers an anomaly detector from developed markets
to a frontier market. A further problem concerns evaluation. Many time-series anomaly
detection results are reported with point adjustment, which can make random scores look
competitive~\cite{kim2022towards}, and the size of this effect on financial data is unknown.

As a result, there is no tested method for flagging abnormal trading days on the DSE that
accounts for changing volatility, and no reliable estimate of how well such a method would
perform.

\section{Research Objectives}

\subsection*{General Objective}
To develop and evaluate a volatility-aware deep anomaly detector for daily stock returns
that can be transferred from global market indices to the Dhaka Stock Exchange.

\subsection*{Specific Objectives}
\begin{enumerate}
    \item To build a leakage-safe data pipeline for global index and DSE daily data,
    including corrections for DSE floor prices, bonus shares and price limits.
    \item To fit GARCH, EGARCH and GJR models on training data and select one on
    out-of-sample residual diagnostics.
    \item To construct anomaly labels from known events, extreme standardized residuals and
    injected anomalies, with the rules fixed before any detector is trained.
    \item To compare an LSTM autoencoder trained on raw returns with the same model trained
    on GARCH-standardized returns, against statistical and classical baselines.
    \item To measure how well detectors trained on global indices transfer to the DSE, with
    and without fine-tuning.
    \item To report event-based and point-adjusted scores over several random seeds, with an
    analysis by volatility regime.
\end{enumerate}

\section{Research Questions and Hypotheses}
\begin{enumerate}
    \item[\textbf{RQ1.}] Does standardizing returns by GARCH-family conditional volatility
    improve deep anomaly detection on stock returns?

    \textit{H1:} An autoencoder trained on GARCH-standardized returns achieves a higher
    event-based F1 and raises fewer false alarms in high-volatility periods than the same
    autoencoder trained on raw returns.

    \item[\textbf{RQ2.}] Can an anomaly detector trained on global market indices be
    transferred to the DSE?

    \textit{H2:} A detector trained on global indices and applied to DSE data without
    retraining outperforms the statistical baselines on the DSE, and fine-tuning on DSE
    training data improves it further.

    \item[\textbf{RQ3.}] How much does point adjustment inflate reported anomaly detection
    performance on financial data?

    \textit{H3:} Point-adjusted F1 is substantially higher than event-based F1 for every
    detector, including the simple baselines.
\end{enumerate}

\section{Scope and Delimitation}
The study uses daily data only, because free intraday DSE data do not exist. On the DSE it
covers the DSEX index from its launch in January 2013 and 10 to 20 liquid stocks chosen by
turnover. Periods under BSEC floor prices are treated as a separate regime. Global market
indices, starting with the S\&P~500, serve as the source domain for transfer. Market
manipulation cases, order-book data, and news or text data are outside the scope.

\section{Significance of the Study}
For the DSE, the study would provide the first out-of-sample test of deep anomaly detection
among the work reviewed, together with documented handling of the exchange's data problems.
For anomaly detection on financial data more generally, it tests whether removing the
volatility level from the input changes which days a deep detector flags, and it measures
how much point adjustment overstates performance. The code will be public, and the labelled
DSE series will be released if the data licence allows, which would give later work a
financial benchmark that current collections lack~\cite{paparrizos2022tsb,schmidl2022anomaly}.

\section{Definition of Terms}
\begin{itemize}
    \item \textbf{Anomaly:} a trading day whose return is unusual relative to the
    volatility of its own period, or that coincides with a documented market event.
    \item \textbf{Volatility clustering:} the tendency of large price changes to follow
    large changes and small changes to follow small ones.
    \item \textbf{GARCH-standardized return:} the return innovation divided by its
    conditional standard deviation, $z_t = \varepsilon_t / \sigma_t$.
    \item \textbf{Transfer learning:} training a model on one domain, here global market
    indices, and applying it, with or without further training, to another, here the DSE.
    \item \textbf{Point adjustment:} an evaluation rule that counts every point of a
    labelled anomalous segment as detected if any one of its points is detected.
\end{itemize}

\section{Organization of the Thesis}
Chapter~2 reviews related work on volatility modelling, financial and time-series anomaly
detection, evaluation practice and transfer learning. Chapter~3 describes the proposed data
pipeline, labels, models and evaluation protocol. Chapter~4 reports preliminary results and
the expected outcomes. Chapter~5 sets out the work plan, risks and fallbacks.
